% GENERATED FILE. Edit canonical lessons, then run npm run generate:books.
% canonical-source-hash: fnv1a64:7682fed4ca44b3e7
% canonical-lessons: LA-C423-semita, LA-C423-callis, LA-C423-trames, LA-C423-clivus, LA-C423-gradior

\chapter{Footpath, Mountain track, Side path, Slope, Step}
\label{ch:la-footpath-mountain-track-side-path-slope-step}
\hlchaptermodality{423}

\begin{chapteropening}
\textbf{By the end of this chapter:} \emph{I can talk about the smaller ways a traveller takes: a footpath, a mountain track, a side path, a slope, and stepping along.}

Complete the last lesson of chapter 423: I can talk about the smaller ways a traveller takes: a footpath, a mountain track, a side path, a slope, and stepping along.
\end{chapteropening}

\section[sēmita, sēmitae]{sēmita, sēmitae — a footpath; a narrow way, a lane}

\label{lesson:LA-C423-semita}



\begin{quote}
Before the new one: say the Latin for a ditch, then the Latin for a rampart.
\end{quote}

\subsection*{You'll want to know: sēmita, sēmitae}
\textbf{sēmita, sēmitae} — \textquotedblleft{}a footpath; a narrow way, a lane\textquotedblright{}.

\begin{quote}
Viātor nōn viā sed sēmitā per silvam ambulat.
\end{quote}

\begin{quote}
\emph{The traveller walks through the wood not on the road but on a footpath.}
\end{quote}

First declension, feminine: \emph{sēmita, sēmitae f.} Narrower than a \emph{via}: room for feet, not for carts.

\begin{cousinweb}
\emph{sēmita} is inherited: EtymDB traces it to Proto-Italic \emph{sēmitā} and lists \emph{trāmes} 'by-path' as a cousin. Lewis \& Short explain it as a way that goes 'aside' (\emph{sē-}) from the main road.
\end{cousinweb}

\begin{grammarlens}[title={What you've built}]
One more word for telling, explaining and describing at B1.
\end{grammarlens}

\subsection*{Your turn}
\begin{itemize}
\raggedright
  \item \emph{Say it:} \emph{sēmita, sēmitae}
  \item \emph{Say it:} \emph{sēmita, sēmitae}, once more
  \item \emph{Say it:} a sentence of your own with \emph{sēmita, sēmitae}
  \item \emph{Recall:} say the Latin for a ditch, then the Latin for a rampart, then say \emph{sēmita, sēmitae} again
\end{itemize}

\begin{tcolorbox}[breakable,colback=teal!4,colframe=teal!35!black,title={Before you move on}]
What does sēmita, sēmitae mean? (\textquotedblleft{}A footpath; a narrow way, a lane\textquotedblright{}.) Say it once more.
\end{tcolorbox}
\section[callis, callis]{callis, callis — a mountain track; a rough path made by cattle}

\label{lesson:LA-C423-callis}



\begin{quote}
Before the new one: say the Latin for a rampart, then the Latin for a footpath.
\end{quote}

\subsection*{You'll want to know: callis, callis}
\textbf{callis, callis} — \textquotedblleft{}a mountain track; a rough path made by cattle\textquotedblright{}.

\begin{quote}
Pastor gregem per callem in montem dūcit.
\end{quote}

\begin{quote}
\emph{The shepherd leads his flock up the mountain along a rough track.}
\end{quote}

Third declension, usually masculine: \emph{callis, callis m.} (sometimes feminine).

\begin{cousinweb}
\emph{callis} is inherited. EtymDB connects it with PIE \emph{kel-} 'to drive, to set in motion' and with \emph{celer} 'quick'; Lewis \& Short give no etymology, so the link stays a likely one rather than a sure one.
\end{cousinweb}

\begin{grammarlens}[title={What you've built}]
One more word for telling, explaining and describing at B1.
\end{grammarlens}

\subsection*{Your turn}
\begin{itemize}
\raggedright
  \item \emph{Say it:} \emph{callis, callis}
  \item \emph{Say it:} \emph{callis, callis}, once more
  \item \emph{Say it:} a sentence of your own with \emph{callis, callis}
  \item \emph{Recall:} say the Latin for a rampart, then the Latin for a footpath, then say \emph{callis, callis} again
\end{itemize}

\begin{tcolorbox}[breakable,colback=teal!4,colframe=teal!35!black,title={Before you move on}]
What does callis, callis mean? (\textquotedblleft{}A mountain track; a rough path made by cattle\textquotedblright{}.) Say it once more.
\end{tcolorbox}
\section[trāmes, trāmitis]{trāmes, trāmitis — a side path, a by-way; a cross path}

\label{lesson:LA-C423-trames}



\begin{quote}
Before the new one: say the Latin for a footpath, then the Latin for a mountain track.
\end{quote}

\subsection*{You'll want to know: trāmes, trāmitis}
\textbf{trāmes, trāmitis} — \textquotedblleft{}a side path, a by-way; a cross path\textquotedblright{}.

\begin{quote}
Mercātor trāmitem scit et ante aliōs ad oppidum pervenit.
\end{quote}

\begin{quote}
\emph{The merchant knows a side path and reaches the town before the others.}
\end{quote}

Third declension, masculine: \emph{trāmes, trāmitis m.} A short cut across country, off the high road.

\begin{cousinweb}
\emph{trāmes} is inherited. Lewis \& Short link it with \emph{trāns} 'across', so it is a path that cuts across country, and EtymDB lists \emph{sēmita} as a cousin.
\end{cousinweb}

\begin{grammarlens}[title={What you've built}]
One more word for telling, explaining and describing at B1.
\end{grammarlens}

\subsection*{Your turn}
\begin{itemize}
\raggedright
  \item \emph{Say it:} \emph{trāmes, trāmitis}
  \item \emph{Say it:} \emph{trāmes, trāmitis}, once more
  \item \emph{Say it:} a sentence of your own with \emph{trāmes, trāmitis}
  \item \emph{Recall:} say the Latin for a footpath, then the Latin for a mountain track, then say \emph{trāmes, trāmitis} again
\end{itemize}

\begin{tcolorbox}[breakable,colback=teal!4,colframe=teal!35!black,title={Before you move on}]
What does trāmes, trāmitis mean? (\textquotedblleft{}A side path, a by-way; a cross path\textquotedblright{}.) Say it once more.
\end{tcolorbox}
\section[clīvus, clīvī]{clīvus, clīvī — a slope; a sloping street, an uphill road}

\label{lesson:LA-C423-clivus}



\begin{quote}
Before the new one: say the Latin for a mountain track, then the Latin for a side path.
\end{quote}

\subsection*{You'll want to know: clīvus, clīvī}
\textbf{clīvus, clīvī} — \textquotedblleft{}a slope; a sloping street, an uphill road\textquotedblright{}.

\begin{quote}
Clīvus longus est, et mūlus lentē ascendit.
\end{quote}

\begin{quote}
\emph{The slope is long, and the mule climbs slowly.}
\end{quote}

Second declension, masculine: \emph{clīvus, clīvī m.} A street that climbs a hill is also a \emph{clīvus}. English \emph{declivity} and \emph{proclivity} come from the same root.

\begin{cousinweb}
\emph{clīvus} is inherited from PIE \emph{\'{k}leywós}, built on \emph{\'{k}ley-} 'to lean, to incline' (EtymDB). English \emph{declivity} and \emph{proclivity} come from the related Latin adjectives \emph{dēclīvis} and \emph{prōclīvis} (Lewis \& Short).
\end{cousinweb}

\begin{grammarlens}[title={What you've built}]
One more word for telling, explaining and describing at B1.
\end{grammarlens}

\subsection*{Your turn}
\begin{itemize}
\raggedright
  \item \emph{Say it:} \emph{clīvus, clīvī}
  \item \emph{Say it:} \emph{clīvus, clīvī}, once more
  \item \emph{Say it:} a sentence of your own with \emph{clīvus, clīvī}
  \item \emph{Recall:} say the Latin for a mountain track, then the Latin for a side path, then say \emph{clīvus, clīvī} again
\end{itemize}

\begin{tcolorbox}[breakable,colback=teal!4,colframe=teal!35!black,title={Before you move on}]
What does clīvus, clīvī mean? (\textquotedblleft{}A slope; a sloping street, an uphill road\textquotedblright{}.) Say it once more.
\end{tcolorbox}
\section[gradior, gradī]{gradior, gradī — to step, to walk}

\label{lesson:LA-C423-gradior}



\begin{quote}
Before the new one: say the Latin for a side path, then the Latin for a slope.
\end{quote}

\subsection*{You'll want to know: gradior, gradī}
\textbf{gradior, gradī} — \textquotedblleft{}to step, to walk\textquotedblright{}.

\begin{quote}
In viā angustā senex lentē graditur.
\end{quote}

\begin{quote}
\emph{On the narrow road the old man walks slowly.}
\end{quote}

Deponent: \emph{gradior, gradī, gressus sum}. Passive endings, active meaning. Its compounds give English \emph{ingredient}, \emph{progress} and \emph{egress}.

\begin{cousinweb}
\emph{gradior} is inherited from PIE \emph{g\textsuperscript{h}red\textsuperscript{h}-} 'to walk, to go' (EtymDB), and \emph{gradus} 'a step' is built on the same root. English \emph{gradient} and \emph{grade} go back to this family (Lewis \& Short).
\end{cousinweb}

\begin{grammarlens}[title={What you've built}]
One more word for telling, explaining and describing at B1.
\end{grammarlens}

\subsection*{Your turn}
\begin{itemize}
\raggedright
  \item \emph{Say it:} \emph{gradior, gradī}
  \item \emph{Say it:} \emph{gradior, gradī}, once more
  \item \emph{Say it:} a sentence of your own with \emph{gradior, gradī}
  \item \emph{Recall:} say the Latin for a side path, then the Latin for a slope, then say \emph{gradior, gradī} again
\end{itemize}

\begin{tcolorbox}[breakable,colback=teal!4,colframe=teal!35!black,title={Before you move on}]
What does gradior, gradī mean? (\textquotedblleft{}To step, to walk\textquotedblright{}.) Say it once more.
\end{tcolorbox}
